\documentclass[a4paper,12pt]{article}

\usepackage[T2A]{fontenc}
\usepackage[utf8]{inputenc}
\usepackage[russian]{babel}

\usepackage{amsmath, amssymb, amsthm, mathtools}
\usepackage{helvet}
\renewcommand{\familydefault}{\sfdefault}
\usepackage{icomma}
\usepackage{geometry}
\usepackage{microtype}
\usepackage{tikz}
\usepackage{tkz-euclide}
\usepackage{pgfplots}
\pgfplotsset{compat=1.18}
\usepackage{tabularx, booktabs, longtable}
\usepackage{enumitem}
\usepackage{hyperref}
\usepackage{parskip}
\usepackage{fancyhdr}
\usepackage{setspace}
\usepackage{xcolor}

\geometry{left=18mm,right=18mm,top=18mm,bottom=20mm}
\onehalfspacing

\setlist[itemize]{leftmargin=1.2em,itemsep=0.25em,topsep=0.25em}
\setlist[enumerate]{leftmargin=1.4em,itemsep=0.35em,topsep=0.35em}

\definecolor{accent}{HTML}{008080}
\definecolor{darkaccent}{HTML}{004D4D}
\definecolor{textgray}{HTML}{333333}
\definecolor{softgray}{HTML}{F2F4F4}

\hypersetup{
    colorlinks=true,
    linkcolor=darkaccent,
    urlcolor=darkaccent
}

\pagestyle{fancy}
\fancyhf{}
\lhead{\textcolor{textgray}{Тригонометрические уравнения}}
\rhead{\textcolor{textgray}{\thepage}}
\renewcommand{\headrulewidth}{0.4pt}
\renewcommand{\headrule}{\hbox to\headwidth{\color{accent}\leaders\hrule height \headrulewidth\hfill}}

\newcommand{\sectionline}{\vspace{0.2em}\textcolor{accent}{\hrule height 0.8pt}\vspace{0.7em}}
\newcommand{\important}[1]{\textcolor{darkaccent}{\textbf{#1}}}
\newcommand{\Z}{\mathbb{Z}}
\newcommand{\tg}{\mathop{\mathrm{tg}}\nolimits}
\newcommand{\ctg}{\mathop{\mathrm{ctg}}\nolimits}
\newcommand{\arctg}{\mathop{\mathrm{arctg}}\nolimits}
\newcommand{\arcctg}{\mathop{\mathrm{arcctg}}\nolimits}

\tkzSetUpLine[line width=1pt, color=accent]
\tkzSetUpPoint[color=accent, fill=accent, size=3]
\tkzSetUpLabel[color=black, font=\small]

\begin{document}

\begin{titlepage}
\thispagestyle{empty}
\centering

\vspace*{30mm}

{\Huge\bfseries Чек-лист\par}
\vspace{8mm}
{\Large\bfseries Тригонометрические уравнения\par}
\vspace{4mm}
{\large Общие формулы решений, аркфункции, приведение углов и замена переменной\par}

\vspace{22mm}

{\large Лёвин Артём Александрович эксклюзивно для Тимофея Васильченко\par}

\vspace{12mm}

{\large \today\par}

\vfill

\begin{tikzpicture}[remember picture,overlay]
\draw[accent, line width=1.2pt]
    ([xshift=25mm,yshift=22mm]current page.south west)
    .. controls ([xshift=65mm,yshift=38mm]current page.south west)
    and ([xshift=-70mm,yshift=10mm]current page.south east)
    .. ([xshift=-25mm,yshift=28mm]current page.south east);
\draw[darkaccent, line width=0.6pt]
    ([xshift=38mm,yshift=15mm]current page.south west)
    .. controls ([xshift=90mm,yshift=28mm]current page.south west)
    and ([xshift=-95mm,yshift=40mm]current page.south east)
    .. ([xshift=-38mm,yshift=18mm]current page.south east);
\end{tikzpicture}

\end{titlepage}

\section*{1. Главная идея занятия}
\sectionline

При решении тригонометрических уравнений важно добиться двух вещей:

\begin{enumerate}
    \item \important{углы должны стать одинаковыми;}
    \item \important{тригонометрические функции желательно привести к одной функции.}
\end{enumerate}

Если в уравнении встречаются разные углы, например \(2x\) и \(x+\frac{3\pi}{2}\), то сначала применяют тригонометрические формулы. После этого уравнение часто сводится к алгебраическому уравнению относительно \(\sin x\), \(\cos x\), \(\tg x\) или \(\ctg x\).

\section*{2. Что такое арксинус, арккосинус и арктангенс}
\sectionline

\subsection*{Смысл аркфункций}

Аркфункция помогает записать угол, если значение тригонометрической функции известно.

\begin{itemize}
    \item \(\arcsin a\) --- это угол, синус которого равен \(a\).
    \item \(\arccos a\) --- это угол, косинус которого равен \(a\).
    \item \(\arctg a\) --- это угол, тангенс которого равен \(a\).
\end{itemize}

Например:
\[
\arcsin \frac12=\frac{\pi}{6},
\qquad
\arcsin \frac{\sqrt2}{2}=\frac{\pi}{4},
\qquad
\arcsin \frac{\sqrt3}{2}=\frac{\pi}{3}.
\]

Если значение нетабличное, например
\[
\arccos \frac{1}{\sqrt{10}},
\]
то его оставляют в таком виде.

\subsection*{Области значений аркфункций}

Для корректной записи решений важно помнить стандартные промежутки:

\[
\arcsin a\in\left[-\frac{\pi}{2};\frac{\pi}{2}\right],
\qquad
\arccos a\in[0;\pi],
\qquad
\arctg a\in\left(-\frac{\pi}{2};\frac{\pi}{2}\right).
\]

\section*{3. Общие формулы для простейших уравнений}
\sectionline

\subsection*{Уравнение \(\sin x=a\)}

Если \(-1\le a\le 1\), то:
\[
x=(-1)^n\arcsin a+\pi n,\qquad n\in\Z.
\]

Также можно использовать раздельную запись:
\[
x=\arcsin a+2\pi n
\quad \text{или} \quad
x=\pi-\arcsin a+2\pi n,
\qquad n\in\Z.
\]

Если \(a>0\), это выглядит особенно наглядно:
\[
\sin x=a
\quad\Longrightarrow\quad
\begin{cases}
x=\arcsin a+2\pi n,\\
x=\pi-\arcsin a+2\pi n,
\end{cases}
\qquad n\in\Z.
\]

Если \(a>0\), то:
\[
\sin x=-a
\quad\Longrightarrow\quad
\begin{cases}
x=-\arcsin a+2\pi n,\\
x=\pi+\arcsin a+2\pi n,
\end{cases}
\qquad n\in\Z.
\]

\subsection*{Уравнение \(\cos x=a\)}

Если \(-1\le a\le 1\), то:
\[
x=\pm\arccos a+2\pi n,\qquad n\in\Z.
\]

Если \(a>0\), то:
\[
\cos x=a
\quad\Longrightarrow\quad
x=\pm\arccos a+2\pi n,\qquad n\in\Z.
\]

Если \(a>0\), то:
\[
\cos x=-a
\quad\Longrightarrow\quad
x=\pi\pm\arccos a+2\pi n,\qquad n\in\Z.
\]

\subsection*{Уравнение \(\tg x=a\)}

\[
\tg x=a
\quad\Longrightarrow\quad
x=\arctg a+\pi n,\qquad n\in\Z.
\]

Если \(a>0\), то:
\[
\tg x=-a
\quad\Longrightarrow\quad
x=-\arctg a+\pi n,\qquad n\in\Z.
\]

\subsection*{Уравнение \(\ctg x=a\)}

\[
\ctg x=a
\quad\Longrightarrow\quad
x=\arcctg a+\pi n,\qquad n\in\Z.
\]

Для ЕГЭ часто удобнее сводить уравнения с котангенсом к тангенсу:
\[
\ctg x=a
\quad\Longleftrightarrow\quad
\tg x=\frac{1}{a},
\qquad a\ne 0.
\]

\section*{4. Табличные значения}
\sectionline

\begin{center}
\renewcommand{\arraystretch}{1.45}
\begin{tabularx}{\linewidth}{@{}cXXXX@{}}
\toprule
\textbf{Угол} & \(0\) & \(\frac{\pi}{6}\) & \(\frac{\pi}{4}\) & \(\frac{\pi}{3}\) \\
\midrule
\(\sin x\) & \(0\) & \(\frac12\) & \(\frac{\sqrt2}{2}\) & \(\frac{\sqrt3}{2}\) \\
\(\cos x\) & \(1\) & \(\frac{\sqrt3}{2}\) & \(\frac{\sqrt2}{2}\) & \(\frac12\) \\
\(\tg x\) & \(0\) & \(\frac{\sqrt3}{3}\) & \(1\) & \(\sqrt3\) \\
\(\ctg x\) & не определён & \(\sqrt3\) & \(1\) & \(\frac{\sqrt3}{3}\) \\
\bottomrule
\end{tabularx}
\end{center}

Дополнительно:
\[
\sin\frac{\pi}{2}=1,
\qquad
\cos\frac{\pi}{2}=0,
\qquad
\tg\frac{\pi}{2}\ \text{не определён},
\qquad
\ctg\frac{\pi}{2}=0.
\]

\section*{5. Простейшие примеры}
\sectionline

\subsection*{Пример 1}

Решить уравнение:
\[
\sin x=\frac{\sqrt3}{2}.
\]

Так как
\[
\arcsin\frac{\sqrt3}{2}=\frac{\pi}{3},
\]
получаем:
\[
x=\frac{\pi}{3}+2\pi n
\quad \text{или} \quad
x=\pi-\frac{\pi}{3}+2\pi n.
\]
Ответ:
\[
x=\frac{\pi}{3}+2\pi n
\quad \text{или} \quad
x=\frac{2\pi}{3}+2\pi n,
\qquad n\in\Z.
\]

\subsection*{Пример 2}

Решить уравнение:
\[
\cos x=-\frac{\sqrt3}{2}.
\]

Так как
\[
\arccos\frac{\sqrt3}{2}=\frac{\pi}{6},
\]
то:
\[
x=\pi\pm\frac{\pi}{6}+2\pi n.
\]
Ответ:
\[
x=\frac{5\pi}{6}+2\pi n
\quad \text{или} \quad
x=\frac{7\pi}{6}+2\pi n,
\qquad n\in\Z.
\]

\section*{6. Формулы, которые чаще всего нужны для усложнения уравнений}
\sectionline

\subsection*{Формулы сложения}

\[
\cos(\alpha+\beta)=\cos\alpha\cos\beta-\sin\alpha\sin\beta.
\]

\[
\sin(\alpha+\beta)=\sin\alpha\cos\beta+\cos\alpha\sin\beta.
\]

\[
\cos(\alpha-\beta)=\cos\alpha\cos\beta+\sin\alpha\sin\beta.
\]

\[
\sin(\alpha-\beta)=\sin\alpha\cos\beta-\cos\alpha\sin\beta.
\]

\subsection*{Формулы двойного угла}

\[
\sin 2x=2\sin x\cos x.
\]

\[
\cos 2x=\cos^2x-\sin^2x.
\]

\[
\cos 2x=1-2\sin^2x.
\]

\[
\cos 2x=2\cos^2x-1.
\]

Выбор формы зависит от того, какая функция уже есть в уравнении. Если рядом стоит \(\sin x\), часто удобно использовать:
\[
\cos 2x=1-2\sin^2x.
\]
Если рядом стоит \(\cos x\), часто удобно использовать:
\[
\cos 2x=2\cos^2x-1.
\]

\subsection*{Основное тригонометрическое тождество}

\[
\sin^2x+\cos^2x=1.
\]

Отсюда:
\[
\sin^2x=1-\cos^2x,
\qquad
\cos^2x=1-\sin^2x.
\]

\section*{7. Алгоритм решения более сложного тригонометрического уравнения}
\sectionline

\begin{enumerate}
    \item Посмотрите, одинаковые ли углы в уравнении.
    \item Если углы разные, примените формулы сложения, приведения или двойного угла.
    \item Постарайтесь привести всё к одной функции: только к \(\sin x\), только к \(\cos x\), только к \(\tg x\) и т.д.
    \item Если появляется повторяющийся элемент, например \(\sin x\), сделайте замену:
    \[
    t=\sin x.
    \]
    \item Обязательно запишите ограничение:
    \[
    -1\le t\le 1
    \quad \text{для } t=\sin x \text{ или } t=\cos x.
    \]
    \item Решите полученное алгебраическое уравнение.
    \item Отбросьте значения \(t\), которые не попадают в допустимый промежуток.
    \item Сделайте обратную замену и решите простейшее тригонометрическое уравнение.
\end{enumerate}

\section*{8. Пример с заменой переменной}
\sectionline

Решить уравнение:
\[
\cos 2x-\cos\left(\frac{3\pi}{2}+x\right)+2=0.
\]

Сначала раскроем второй косинус:
\[
\cos\left(\frac{3\pi}{2}+x\right)
=
\cos\frac{3\pi}{2}\cos x-\sin\frac{3\pi}{2}\sin x.
\]
Так как
\[
\cos\frac{3\pi}{2}=0,
\qquad
\sin\frac{3\pi}{2}=-1,
\]
получаем:
\[
\cos\left(\frac{3\pi}{2}+x\right)=\sin x.
\]

Тогда исходное уравнение принимает вид:
\[
\cos2x-\sin x+2=0.
\]

Используем формулу:
\[
\cos2x=1-2\sin^2x.
\]
Получаем:
\[
1-2\sin^2x-\sin x+2=0.
\]
\[
-2\sin^2x-\sin x+3=0.
\]
Умножим на \(-1\):
\[
2\sin^2x+\sin x-3=0.
\]

Сделаем замену:
\[
t=\sin x,
\qquad
-1\le t\le 1.
\]
Тогда:
\[
2t^2+t-3=0.
\]

Дискриминант:
\[
D=1-4\cdot2\cdot(-3)=25.
\]

Корни:
\[
t_1=\frac{-1+5}{4}=1,
\qquad
t_2=\frac{-1-5}{4}=-\frac32.
\]

Корень \(t_2=-\frac32\) не подходит, так как
\[
-\frac32<-1.
\]

Остаётся:
\[
\sin x=1.
\]
Ответ:
\[
x=\frac{\pi}{2}+2\pi n,
\qquad n\in\Z.
\]

\section*{9. Как проверять корни после замены}
\sectionline

Если сделана замена
\[
t=\sin x
\quad \text{или} \quad
t=\cos x,
\]
то обязательно:
\[
-1\le t\le 1.
\]

Например, если получен корень
\[
t=\frac{7+\sqrt{33}}{4},
\]
то он не подходит, так как
\[
\frac{7+\sqrt{33}}{4}>1.
\]
Действительно:
\[
\frac{7+\sqrt{33}}{4}>1
\quad\Longleftrightarrow\quad
7+\sqrt{33}>4
\quad\Longleftrightarrow\quad
\sqrt{33}>-3,
\]
что верно. Значит, число больше \(1\).

Если получен корень
\[
t=\frac{7-\sqrt{33}}{4},
\]
то нужно проверить, что он лежит в промежутке \([-1;1]\). Так как \(\sqrt{33}<7\), число положительно. Достаточно проверить, что оно меньше \(1\):
\[
\frac{7-\sqrt{33}}{4}<1
\quad\Longleftrightarrow\quad
7-\sqrt{33}<4
\quad\Longleftrightarrow\quad
3<\sqrt{33}
\quad\Longleftrightarrow\quad
9<33.
\]
Значит,
\[
0<\frac{7-\sqrt{33}}{4}<1.
\]

\section*{10. Типичные ошибки}
\sectionline

\begin{enumerate}
    \item \important{Пишут \(2\pi n\) вместо \(\pi n\) для тангенса.}  
    Для \(\tg x=a\) период равен \(\pi\), поэтому:
    \[
    x=\arctg a+\pi n.
    \]

    \item \important{Забывают второй корень у синуса или косинуса.}  
    Уравнения \(\sin x=a\) и \(\cos x=a\) обычно дают две серии решений на периоде \(2\pi\).

    \item \important{Не проверяют ограничение после замены.}  
    Если \(t=\sin x\) или \(t=\cos x\), то корень \(t\) обязан лежать в \([-1;1]\).

    \item \important{Меняют знак неравенства неправильно.}  
    При умножении или делении неравенства на отрицательное число знак меняется на противоположный.

    \item \important{Не приводят углы к одинаковому виду.}  
    Если в уравнении есть \(2x\), \(x+\frac{3\pi}{2}\), \(x\), сначала нужно использовать формулы.

    \item \important{Используют не ту формулу двойного угла.}  
    Формула выбирается под нужную функцию. Если требуется перейти к \(\sin x\), используйте
    \[
    \cos2x=1-2\sin^2x.
    \]

    \item \important{Теряют условие \(n\in\Z\).}  
    В ответе к тригонометрическому уравнению обязательно указывать:
    \[
    n\in\Z.
    \]
\end{enumerate}

\section*{11. Короткая памятка}
\sectionline

\begin{itemize}
    \item \(\sin x=a\): две серии решений, если \(-1<a<1\).
    \item \(\cos x=a\): удобно писать \(x=\pm\arccos a+2\pi n\).
    \item \(\tg x=a\): период \(\pi\), а не \(2\pi\).
    \item Разные углы нужно привести к одинаковым через формулы.
    \item После замены \(t=\sin x\) или \(t=\cos x\) проверяйте \(-1\le t\le1\).
    \item Если значение табличное, аркфункцию можно заменить конкретным углом.
    \item Если значение нетабличное, аркфункцию оставляют в ответе.
\end{itemize}

\newpage

\section*{Тренировочный блок}
\sectionline

\textbf{Решите уравнения.} Ответы запишите с указанием \(n\in\Z\). Задачи расположены по нарастающей сложности.

\begin{enumerate}
    \item Решите уравнение:
    \[
    \sin x=\frac{\sqrt3}{2}.
    \]

    \item Решите уравнение:
    \[
    \cos x=-\frac{\sqrt2}{2}.
    \]

    \item Решите уравнение:
    \[
    \tg x=-\sqrt3.
    \]

    \item Решите уравнение:
    \[
    \sin x=-\frac12.
    \]

    \item Решите уравнение:
    \[
    \cos x=\frac{1}{\sqrt{10}}.
    \]

    \item Решите уравнение:
    \[
    \cos 2x-\sin x+2=0.
    \]

    \item Решите уравнение:
    \[
    \cos 2x+\cos x-2=0.
    \]

    \item Решите уравнение:
    \[
    \cos\left(\frac{3\pi}{2}+x\right)+2\sin^2x-3\sin x=0.
    \]

    \item Решите уравнение:
    \[
    2\cos^2x-3\cos x+1=0.
    \]

    \item Решите уравнение:
    \[
    \sin\left(x+\frac{\pi}{2}\right)+\cos2x=0.
    \]
\end{enumerate}

\newpage

\section*{Ответы}
\sectionline

\begin{center}
\renewcommand{\arraystretch}{1.55}
\begin{tabularx}{\linewidth}{@{}cX@{}}
\toprule
\textbf{№} & \textbf{Ответ} \\
\midrule
1 &
\[
x=\frac{\pi}{3}+2\pi n
\quad \text{или} \quad
x=\frac{2\pi}{3}+2\pi n,\quad n\in\Z.
\]
\\

2 &
\[
x=\frac{3\pi}{4}+2\pi n
\quad \text{или} \quad
x=\frac{5\pi}{4}+2\pi n,\quad n\in\Z.
\]
\\

3 &
\[
x=-\frac{\pi}{3}+\pi n,\quad n\in\Z.
\]
\\

4 &
\[
x=-\frac{\pi}{6}+2\pi n
\quad \text{или} \quad
x=\frac{7\pi}{6}+2\pi n,\quad n\in\Z.
\]
\\

5 &
\[
x=\pm\arccos\frac{1}{\sqrt{10}}+2\pi n,\quad n\in\Z.
\]
\\

6 &
\[
x=\frac{\pi}{2}+2\pi n,\quad n\in\Z.
\]
\\

7 &
\[
x=2\pi n,\quad n\in\Z.
\]
\\

8 &
\[
x=\frac{\pi}{2}+2\pi n
\quad \text{или} \quad
x=\frac{7\pi}{6}+2\pi n
\quad \text{или} \quad
x=\frac{11\pi}{6}+2\pi n,\quad n\in\Z.
\]
\\

9 &
\[
x=2\pi n
\quad \text{или} \quad
x=\pm\frac{\pi}{3}+2\pi n,\quad n\in\Z.
\]
\\

10 &
\[
x=\frac{\pi}{3}+2\pi n
\quad \text{или} \quad
x=\frac{5\pi}{3}+2\pi n,\quad n\in\Z.
\]
\\
\bottomrule
\end{tabularx}
\end{center}

\end{document}